\documentclass[a4paper,10pt]{elsarticle}
\usepackage{lineno,hyperref}
\usepackage{booktabs}
\usepackage{listings}
\usepackage{xcolor}
\usepackage{graphicx}
\usepackage{float}
\usepackage{longtable}
\usepackage{amsmath}
\begin{document}

\subsection*{Abstract}

Rationale: Imaging data driven artificial intelligence  (AI) research in Nuclear Medicine (NM) increasingly depends on large-scale multi-modal 3D/4D datasets (PET/SPECT with CT/MRI) that must be resampled, aligned and quantitatively normalized without losing the spatial and clinical metadata required for physical interpretation (e.g. SUV, absorbed dose). Existing pipelines often struggle with (1) biobank-scale volume and I/O overhead, (2) execution speed and the ``two-language'' barrier, (3) metadata drift when images are converted to raw arrays (4) and end-to-end differentiability for novel physics-informed machine learning.

\textbf{Methods:} We developed \texttt{MedImages.jl}, a native Julia library that binds voxel data to spatial metadata (origin, spacing, direction) and clinical/temporal fields using a typed \texttt{MedImage} structure and a batched multimodal container (\texttt{BatchedMedImage}) for synchronized transforms. Volume-scale performance was evaluated on a 100-subject autoPET PET/CT alignment pipeline (target grid $512^3$) comparing HDF5-based persistence to a Python caching pipeline. The execution speed was benchmarked on GPU for resampling, orientation changes, and fused affine transforms, and compared with SimpleITK (ITK-based C++ reference toolkit). Differentiability was validated via gradient tests and learned spatial pipelines (inverse-rotation spatial transformer; organ-specific affine registration) and extended to physics-informed neural network dosimetry using a Universal Differential Equation (UDE) for $^{177}$Lu-PSMA trained against Monte Carlo dose ground truths. Metadata integrity was assessed by SUV agreement with SimpleITK and by an end-to-end PET/CT transformation chain with lesion-based readouts.

\textbf{Results:} \texttt{MedImages.jl} achieved a $7.2\times$ turnaround speedup versus caching in the 100-subject benchmark and up to $115\times$ (resampling), $71\times$ (orientation), and $135\times$ (fused affine) GPU speedups versus CPU. SUVbw matched SimpleITK to $10^{-14}$ (double precision), and compounded transforms preserved Mean SUV ($3.4913 \rightarrow 3.4660$; 0.73\%) and lesion volume ($52.55 \rightarrow 51.86$~cm$^3$; 1.32\%), consistent with interpolation effects rather than metadata deterioration. Learned spatial transforms reduced reconstruction error by $>65$\%, and UDE dosimetry refinement converged to MSE $4.14\times10^{-4}$ versus Monte Carlo ground-truth dose maps.

\textbf{Conclusion:} \texttt{MedImages.jl} provides a unified, GPU-accelerated, metadata-aware, and differentiable foundation that addresses scalability, speed, differentiability, and metadata integrity for quantitative nuclear medicine preprocessing and AI/SciML workflows.

\textbf{Key words:} nuclear medicine; PET/CT; SUV; metadata; differentiable programming

\section{Introduction}

Nuclear medicine diagnostic data, such as PET and SPECT, provide quantitative measurements of biologic processes that support oncologic staging, assessment of therapy response, and planning of theranostic treatment. Correct processing and interpretation of functional imaging data require that measured uptake is assigned to the same patient-space coordinates as the anatomic reference (CT or MRI). Hybrid imaging therefore combines complementary modalities within a single, physically defined sampling space specified by image origin, grid or voxel spacing, and orientation. 
 
As multimodal imaging is increasingly reused as input for AI models, managing data volume and maintaining heterogeneous acquisition metadata has become a central practical challenge \cite{}. A high computational data preparation process typically requires: loading of a large number of examinations from disk, mapping each modality into a shared patient-space coordinate system, resampling onto standardized grids and applying geometric transforms for augmentation or registration. Preserving spatial and quantitative metadata is necessary throughout the multistep preprocessing phase so that voxel values remain physically interpretable. 

Metadata management is not only a practical but also a safety-relevant issue. Converting medical image objects (DICOM, NIfTI) to raw arrays for deep learning can decouple voxel intensities from their patient-space definition and from clinical metadata used for quantification. This ``metadata drift'' can yield visually plausible results that are spatially inconsistent or quantitatively invalid \cite{}.

In widely used image-processing environments such as SimpleITK \cite{LowekampChen2013} and MONAI \cite{CardosoLi2022}, data handling is often built around workflows and chains of separate libraries rather than architectures tailored to biobank-scale diagnostic imaging. As a result, a substantial fraction of total processing time is used not only for core preprocessing operations itself, but also to repeatedly copy very large image volumes between intermediate storage steps and different software components, and to ensure the consistency of spatial information \cite{Dexl_2025}.

. 

These limitations become even more pronounced in emerging Scientific Machine Learning (SciML) workflows, which increasingly rely on end-to-end optimization in which geometric preprocessing participates in training of AI models. In nuclear medicine, such workflows include learned registration and motion correction, data augmentation, and physics-informed models for image reconstruction or dosimetry. Although Python-based frameworks provide differentiable operations, practical composability across heterogeneous software components often requires committing the workflow to a single tensor ecosystem, creating friction when integrating third-party tools or classical numerical solvers \cite{PalBhattacharya2024, EschleGl2023}. 

To address these limitations, we developed MedImages.jl, a Julia-based imaging ecosystem that converts high-level code into efficient machine instructions using just-in-time intermediate optimization and enables native implementations of geometric image operations. As a result, these operations can run efficiently on both CPUs and GPUs  \cite{BezansonChen2018}. 

MedImages.jl binds voxel data to spatial and acquisition metadata through typed image representations, so that the typical for preprocessing geometric operations update the underlying physical description rather than operating on detached arrays. Core operators like resampling, affine warping, orientation handling, cropping and padding are implemented as native Julia kernels with transparent GPU acceleration, reducing practical bottlenecks in repeated 3D processing, supporting multimodal synchronization through batched representations.

The new framework is designed for interoperability with the broader Julia SciML ecosystem and with automatic differentiation libraries such as Zygote.jl and Enzyme.jl, which allows for differentiable spatial pipelines \cite{10.1145/3458817.3476165}. This permits geometric image operations to be incorporated directly into end-to-end optimization, allowing spatial transformations to influence model training rather than remain external preprocessing steps. Julia further reduces software fragmentation by enabling image transformations, reconstruction algorithms, numerical solvers, and optimization routines to operate within a single coherent framework instead of across separate software layers that must be manually linked \cite{RoeschGreener2023}. Since these compiler-level differentiation engines can act on arbitrary Julia code, gradients can propagate through solvers and GPU kernels without rewriting the underlying workflow in a specialized tensor framework \cite{Fischer\_2019}. 

In this work, we evaluate MedImages.jl across nuclear medicine–relevant benchmarks and experiments, both in isolation and in integration with complementary Julia libraries, with emphasis on (i) scalability for large PET/CT cohorts with high preprocessing demands, (ii) execution speed without the maintenance costs of a two-language architecture, (iii)  metadata preservation to maintain the physical interpretability required for quantitative nuclear medicine, and (iv) differentiability of geometric operations for learned spatial transforms.


\section{Methods}

MedImages.jl was developed to address several practical requirements that are highly relevant in medical imaging, including preservation of image geometry and acquisition metadata, reliable handling of clinical image formats, efficient processing of large multimodal datasets, and support for optimization-based learning workflows. 
A central design feature is the \verb|MedImage| data structure, which stores voxel values together with essential spatial information such as image position, voxel spacing, orientation, patient identity, and acquisition time. This avoids a common problem in AI-oriented workflows, where images are converted into generic tensors and thereby lose information needed to interpret them in physical space \cite{}. 
For robust reading and writing of clinical formats such as DICOM and NIfTI, MedImages.jl uses ITK through a wrapper interface, whereas spatial operations such as rotation, translation, and resampling are implemented directly in Julia language. This design combines the reliability of a mature medical imaging toolkit for file handling with the flexibility to use spatial operations in modern learning and optimization pipelines. 
For larger studies, the \verb|BatchedMedImage| structure allows multiple images to be processed together while preserving the spatial properties of each individual volume. In combination with HDF5-based storage and GPU acceleration, this architecture reduces data-loading overhead, improves scalability, and supports workflows in which image transformations can be incorporated directly into model optimization.  

Moreover, compatibility with the broader Julia SciML ecosystem allows MedImages.jl to interoperate with other scientific computing libraries and supports differentiable geometric operations for learned spatial transforms.  

To demonstrate that MedImages.jl addresses the central computational and metadata-related challenges of large-scale medical image analysis, we conducted three proof-of-concept experimental series.  All experiments described are open-source and reproducible via the codebase provided in this repository.

\begin{table}[H]
\centering
\caption{Core functionalities and data structures of MedImages.jl for 3D medical image processing.}
\label{tab:medimages_functionalities}
\renewcommand{\arraystretch}{1.2}
\begin{tabular}{|p{0.24\textwidth}|p{0.68\textwidth}|}
\hline
\textbf{Functionality} & \textbf{Description} \\
\hline
\texttt{MedImage} data structure & Unified medical image representation that stores voxel data together with clinical and vectorized spatial metadata. \\
\hline
\texttt{BatchedMedImage} structure & Batched representation for joint processing of multiple images while preserving the spatial properties of each individual volume. \\
\hline
Imaging-format-compatible I/O & Loading of standard medical image formats, including NIfTI and DICOM, via \texttt{ITKIOWrapper}. \\
\hline
HDF5 support & Fast loading and saving in big-data compatible format, facilitating efficient intermediate storage in repeated processing and deep-learning workflows. \\
\hline
GPU acceleration & Transparent GPU support for spatial transformations through \texttt{KernelAbstractions.jl}, enabling accelerated parallel execution of preprocessing and geometric operations. \\
\hline
Resampling & Native Julia functions for resampling, enabling voxelwise correspondence across modalities. \\
\hline
Spatial transformations & Native Julia functions for manipulation of spatial properties of 3D image volumes, including \texttt{rotate\_mi}, \texttt{translate\_mi}, \texttt{scale\_mi}, \texttt{resample\_to\_spacing}, \texttt{crop\_mi}, and \texttt{pad\_mi}. \\
\hline
Orientation management & Tools for handling, preserving, and converting image orientations across neuroimaging and medical imaging coordinate systems. \\
\hline
Automatic differentiation & Differentiable implementation of spatial transformations with well-defined gradients through \texttt{Zygote.jl} and \texttt{Enzyme.jl}. \\
\hline
\end{tabular}
\end{table}




\textbf{Experimental Series 1: Evaluation of Scalability and Execution Speed}

For large-scale imaging studies, we evaluated and compared MedImages.jl with contemporary frameworks with respect to two key requirements: scalability to large multimodal cohorts and runtime of core spatial preprocessing operations in repeated deep-learning workflows.

To assess scalability and computational overhead, we designed a 100-case preprocessing benchmark based on the autoPET dataset \cite{Dexl_2025}. This setup was intended to reflect a realistic large-cohort scenario in which multiple PET/CT studies must be repeatedly loaded, harmonized and prepared during iterative model training.

The main preprocessing task was multimodal spatial alignment of heterogeneous PET and CT volumes on a common isotropic reference grid of 512 × 512 × 512 voxels. This harmonization is essential for multimodal neural networks, which require strict voxelwise spatial correspondence across input channels. In practical terms, alignment means transforming both modalities into the same physical coordinate system such that a voxel in position (i,j,k) of the PET volume corresponds to the same anatomical location as the voxel in position (i,j,k) of the CT volume.

This benchmark was divided into two phases per subject: first, loading the preprocessed case from cache, and second, applying the remaining computational transformations required to prepare the data for training.

Within the MedImages.jl pipeline, HDF5-based storage was combined with a fused affine preprocessing strategy to reduce computational overhead. Rather than performing orientation normalization, voxel spacing adjustment, and spatial centering as separate steps, these operations were merged into a single affine transformation. This allowed each volume to be resampled in a single step, thus reducing repeated interpolation and minimizing data movement.

MedImages.jl was compared with MONAI using its PersistentDataset mechanism \cite{}, thus assessing not only the computational performance at the framework-level, but also the effect of different data persistence strategies. Performance was evaluated under repeated training conditions, beginning from epoch 2 onward, under the assumption that the data had already undergone initial preprocessing steps and had been stored in cache.

To assess raw computational speed independently of the caching benchmark, we additionally measured isolated spatial operations across multiple scales. These tests included resampling, orientation changes, and affine transformations, which represent some of the most common geometric operations in multimodal medical image preprocessing. Benchmarks were performed on single volume of size 256 × 256 × 128 using both CPU and GPU backends.  SimpleITK was evaluated on CPU, reflecting its typical use for such geometric operations, whereas MedImages.jl was evaluated on both CPU and GPU in order to quantify the effect of hardware acceleration on the same class of tasks.

Both the 100-subject caching pipeline benchmark and the isolated raw computational speed benchmarks (resampling, orientation changes, fused affine transformations) were executed on same hardware configurations: an Intel Core Ultra 9 285K processor for CPU-bound computations and an NVIDIA RTX 3090 GPU for hardware-accelerated processing."


 . 

\textbf{Experimental Series 2: Verification of Metadata Integrity and Quantitative Clinical Fidelity  }

To evaluate the quantitative accuracy of MedImages.jl, Standardized Uptake Value normalized by body weight (SUVbw) was compared against a reference implementation in SimpleITK. SUVbw was computed using same PET image data from autopet dataset and corresponding clinical and temporal metadata in both frameworks. The comparison was performed in double-precision floating point arithmetic, and numerical agreement between both implementations was assessed .  For the fused affine interpolation kernel, the results obtained on the CPU  by reference Python library and GPU were compared directly, and numerical agreement was strictly verified within a tolerance of $10^{-4}$ to guarantee floating-point consistency across backends.


An end-to-end consistency experiment was subsequently conducted on the subset of autoPET dataset to assess the preservation of quantitative fidelity under a sequence of spatial transformations. The processing pipeline consisted of (1) segmentation of lesion masks on the original image data, (2) conversion of PET intensities to SUV, (3) resampling of PET volumes and corresponding lesion masks to the native CT coordinate space, (4) resampling to an isotropic voxel spacing of 2.0 mm, (5) application of a 45° rotation around the z-axis, and (6) translation by 10 voxels along the x-axis. Nearest-neighbor interpolation was applied to segmentation masks, while linear interpolation was used for continuous PET and CT intensity data.
Quantitative evaluation was performed by measuring mean SUV and total lesion volume within the segmentation mask before and after the transformation sequence.

To evaluate the capability of MedImages.jl for handling complex theranostic imaging data, a multi-modal processing experiment was conducted using a dataset consisting of four imaging modalities from a single patient who underwent a radioligand therapy. The data set included: (1) [$^{177}$Lu]Lu-PSMA SPECT with attenuation correction (AC) (2) [$^{177}$Lu]Lu-PSMA SPECT without attenuation correction (NAC), (3) a voxel-wise absorbed dose map, and (4) a CT volume serving as anatomical reference. Each modality was assumed to originate from independent acquisition and reconstruction pipelines and therefore retained its own voxel grid, data type, and spatial metadata, including voxel spacing, image origin, and orientation.
All four modalities were loaded into a single BatchedMedImage data structure. During this step, the original voxel data and modality-specific spatial metadata were preserved without prior resampling or harmonization.Subsequently, a geometric transformation was applied to the entire batch. A spatial rotation of $45^{\circ}$ around the Z-axis (axial plane) was executed simultaneously across all four modalities using a single batched transformation call, followed by a linear translation of 10 voxels along the X-axis to verify spatial origin tracking.

The outcome of the transformation was evaluated qualitatively by visual inspection of corresponding axial slices before and after transformation. Particular attention was given to the preservation of spatial alignment between functional tracer distribution and anatomical structures. The consistency of spatial relationships across modalities after transformation served as the primary validation criterion. 


\textbf{Experimental Series 3: Verification of Differentiability and Learned Spatial Transformations}

To determine whether spatial image operations remained mathematically trainable and could therefore be used reliably in end-to-end learning pipelines, gradient propagation was numerically evaluated across representative geometric transformations implemented in MedImages.jl. Reverse-mode gradients were computed with Zygote.jl for translation, padding, cropping, and fused affine interpolation. For translation and padding, correctness was assessed by verifying that the gradient remained constant throughout the transformed image domain. For cropping, the gradient pattern was examined to confirm preservation within the retained image region and suppression outside the cropped area. For the fused affine interpolation kernel, the gradients obtained on the CPU and GPU were compared directly, and numerical agreement was verified within a tolerance of 10 \^ -4. 

To further assess end-to-end differentiability in a trainable setting, a learned inverse rotation experiment was conducted using synthetic single-case batched image configurations. A synthetic 3D structural line image was generated on a strictly confined $16 \times 16 \times 16$ voxel grid to ensure strong spatial gradient signals. A total of 120 batched sample configurations (100 for training, 20 for testing) were generated by applying random three-axis sequential rotations within an angular range of $\pm 30^{\circ}$. A lightweight 3D Convolutional Neural Network (CNN) combined with a Multi-Layer Perceptron (MLP) head was then trained using gradient descent for 20 epochs to predict the inverse Euler rotation parameters directly from the spatial grids.



\subsection{High-Fidelity Dosimetry via Triple-Branch Universal Differential Equations}

In order to showcase the example of the end to end practical implementation utilizing Medimages.jl Authors had designed and executed the Dosimetry experiment utilizing unique SciML julia library .  

\subsubsection{Model Architecture: The Triple-Branch UDE}
The proposed model evolves the standard dosimetry framework into a simplified neural Universal Differential Equation (UDE) that partitions mechanistic physical laws from learned stochastic scattering.  The core innovation is the expansion to a \textbf{Triple-Branch Input} architecture:
\begin{enumerate}
    \item \textbf{Activity Branch ($A_{total}$)}: Tracks the time-varying radiopharmaceutical concentration derived from the compartmental ODE states ($A_{free} + A_{bound}$).
    \item \textbf{Density Branch ($\rho$)}: Incorporates voxel-specific mass-density derived from CT Hounsfield Units.
    \item \textbf{Gradient Branch ($|\nabla \rho|$)}: A high-pass geometric prior that explicitly signals tissue-bone and tissue-air boundaries, allowing the network to learn asymmetric Compton scattering.
\end{enumerate}

\subsubsection{Numerical Stability: Balanced Unit Training}
To ensure convergence on GPU architectures, we implemented \textbf{Balanced Unit Training}. SPECT activity variables are internally re-scaled by $10^{-6}$ to bring the IVP (Initial Value Problem) into an $O(1)$ range. The physical dose constant ($C_{conv}$) is balanced accordingly, ensuring the integrated output maintains its physical identity in Grays (Gy).

\subsection{Physical Formulation}
The system is defined by the following coupled system of equations:
\begin{equation}
    \frac{d}{dt} \text{Dose}(t) = \text{softplus} \left( \frac{A_{total}(t) \cdot C_{conv}}{\rho \cdot V} + \mathcal{N}_\theta(A_{std}, \rho_{std}, |\nabla \rho|_{std}) \right)
\end{equation}
Where $\mathcal{N}_\theta$ is a 3D Convolutional Neural Network learning the corrective residual to the analytical dose rate. 

The formulation integrates three distinct computational logic blocks:
\begin{enumerate}
    \item \textbf{Analytical Physics Term}: The component $\frac{A_{total}(t) \cdot C_{conv}}{\rho \cdot V}$ represents the local energy deposition. Here, $A_{total}(t)$ is the sum of free and bound compartmental activity, $C_{conv}$ is the dose conversion constant, and the denominator $\rho \cdot V$ calculates the voxel mass (Density $\times$ Volume), ensuring the result maintains the physical unit of Grays (Gy).
    \item \textbf{Neural Network Residual}: The term $\mathcal{N}_\theta(A_{std}, \rho_{std}, |\nabla \rho|_{std})$ acts as a learned correction for stochastic effects such as asymmetric Compton scattering. By utilizing a Triple-Branch input (standardized activity, density, and the density gradient $|\nabla \rho|$), the model explicitly identifies tissue boundaries where classical analytical kernels typically fail.
    \item \textbf{Differentiable Positivity}: The entire rate is wrapped in a \textit{softplus} activation function. This ensures that the generated dose remains non-negative while providing a smooth, continuously differentiable surface required for stable backpropagation through the ODE solver during end-to-end SciML training.
\end{enumerate}



\section{Results }
All results in scalability, computational performance, metadata integrity, quantitative fidelity, and differentiability are summarized in Table~\ref{tab:all_results}.

\subsection*{Evaluation of Scalability and Processing Speed}
In iterative training scenarios, MedImages.jl demonstrated substantially reduced processing times compared to established frameworks, decreasing total turnaround per subject from several hundred milliseconds (MONAI) to well below one tenth of a second. The greatest improvements were observed during the transformation stage, where the fused affine pipeline reduces repeated interpolation and data movement by performing multiple spatial operations in a single computational step. 

GPU acceleration further amplified performance gains across core spatial operations. For medium-sized volumes ($256 \times 256 \times 128$), resampling, orientation changes, and affine transformations achieved speedups exceeding one to two orders of magnitude compared with CPU execution, with additional improvements over standard CPU-based libraries. In large-scale settings ($512^3$ target grid, $n=100$), per-subject processing times remained below 200~ms, confirming suitability for high-throughput and biobank-scale workflows.

\subsection*{Evaluation of Metadata Integrity and Quantitative Fidelity}
Quantitative validation demonstrated very high numerical agreement with the reference standard at machine precision for SUV calculations, with differences on the order of $10^{-14}$. Across compound spatial transformations, mean SUV changed to a negligible extent (from 3.49 to 3.47), and lesion volume remained stable (from approximately 52.6~cm$^3$ to 51.9~cm$^3$), corresponding to relative deviations below 1--2\%. These differences are consistent with expected interpolation effects and do not indicate degradation of clinical metadata.

\subsection*{Evaluation of Differentiability and Learned Transformations}
Gradient propagation across spatial operations followed expected analytical behavior, with identity-preserving transformations maintaining uniform gradients and spatially selective operations producing region-dependent responses. Numerical agreement between CPU and GPU gradient computations was maintained within a tolerance of $10^{-4}$.

In learned transformation experiments, optimization of geometric parameters resulted in a reduction in reconstruction error exceeding 65\% within 20 training iterations, demonstrating effective end-to-end training. Affine registration tasks showed stable convergence using GPU-based gradient optimization, confirming that complex spatial transformations can be integrated directly into differentiable learning pipelines without external implementations.


 
\begin{table}[H]
\centering
\caption{Summary of results across all experimental series.}
\label{tab:all_results}
\renewcommand{\arraystretch}{1.2}
\begin{tabular}{|l|l|c|c|}
\hline
\textbf{Experiment Series} & \textbf{Metric / Task} & \textbf{Result} & \textbf{Reference / Comparison} \\
\hline

\multicolumn{4}{|c|}{\textbf{Series 1: Scalability and Execution Speed}} \\
\hline
Data loading & Load from cache & \textbf{~40 ms} (GPU vs GPU)& MONAI: ~150 ms \\
Computation & Transform (compute) & \textbf{~50 ms} (GPU vs GPU)& MONAI: ~500 ms\\
Overall & Total turnaround & \textbf{~90 ms} (GPU vs GPU)& MONAI: ~650 ms\\
GPU performance & Resampling & \textbf{115$\times$ faster (GPU vs CPU)} & 17$\times$ vs SimpleITK \\
GPU performance & Orientation changes & \textbf{71$\times$ faster (GPU vs CPU)} & 31$\times$ vs SimpleITK \\
GPU performance & Fused affine transform & \textbf{135$\times$ faster (GPU vs CPU)} & 8.1$\times$ vs SimpleITK \\
& & & \\
\hline

\multicolumn{4}{|c|}{\textbf{Series 2: Metadata Integrity and Quantitative Fidelity}} \\
\hline
SUV validation & SUVbw agreement & \textbf{$\approx 10^{-14}$} & Matches SimpleITK (machine precision) \\
Transformation consistency & Mean SUV deviation & \textbf{0.73\%} & Before vs after pipeline \\
Transformation consistency & Lesion volume deviation & \textbf{1.32\%} & Before vs after pipeline \\
\hline
\multicolumn{4}{|c|}{\textbf{Series 3: Differentiability and Learned Transformations}} \\
\hline
Gradient validation & Identity transforms & \textbf{Gradient = 1} & Constant gradient propagation \\
Gradient validation & Cropping & \textbf{1 (inside), 0 (outside)} & Expected spatially selective gradient pattern\\
Gradient validation & CPU vs GPU consistency & \textbf{$< 10^{-4}$ difference} & Numerical consistency between hardware backends\\
Learning experiment & Inverse rotation (CNN) & \textbf{> 65\% MSE reduction} & End-to-end optimization\\
Registration & Affine optimization & \textbf{Converged} & GPU-based gradient optimization \\
 UDE dosimetry& Final voxel-wise dose error & \textbf{MSE = $4.14 \times 10^{-4}$} & Monte Carlo ground-truth dose maps\\
\hline
\end{tabular}
\end{table}

\section{ Disscusion}

As biomedical imaging increasingly embraces complex multimodal, multidimensional, and longitudinal studies, alongside scientific machine learning (SciML), the interoperability and performance of data management frameworks become critical. Our results indicate that MedImages.jl’s metadata handling paradigm, combined with seamless integration into the Julia scientific ecosystem, provides a highly performant and differentiable framework that systematically addresses the three core challenges outlined in the Introduction.  

\subsection*{Addressing Challenges 1 and 2 (Volume and Speed): The Advantage of Native Fused Kernels}
Accurate metadata alignment is a fundamental requirement for contemporary multimodal image analysis. For example, training convolutional neural networks (CNNs) on combined PET/CT data requires multi-channel tensors in which voxel (i,j,k) in the PET dataset corresponds precisely to the same anatomical location as corresponding voxel in the CT dataset. Owing to the large size of 3D voxel-based imaging data, the necessary continuous spatial resampling may, as observed in our experiments, represent the principal computational bottleneck in deep learning workflows. 

Integration of metadata within the \texttt{MedImage} structure directly translates into performance gains observed in our benchmarks. The 7.2$\times$ turnaround advantage observed in high-frequency training loops is primarily attributed to MedImages.jl's integration with HDF5 -- which bypasses the serialization costs inherent in many Python caching mechanisms -- and its single-pass Fused Affine kernel. By mathematically collapsing the entire spatial transformation stack (normalization, resampling, centering, etc.) into a single CUDA execution, MedImages.jl eliminates the memory traffic and synchronization overhead associated with sequential operation queues, providing a highly efficient solution for biobank-scale deployments.

\subsection*{Addressing Challenge 2 (Metadata Management): Comparing Metadata Architectures}
Modern Python frameworks like MONAI \cite{CardosoLi2022} have introduced sophisticated mechanisms to couple metadata with image arrays. MONAI utilizes a MetaTensor structure, appending a metadata dictionary to the PyTorch Tensor. A similar design pattern appears in TorchIO \cite{perezgarcia2021torchio}, which prioritizes convenient multimodal coordination but still builds its core processing around external imaging backends. TorchIO approaches the problem through a hierarchical Subject class system. Although effective for multi-modal synchronization, TorchIO is primarily a manager of third-party backends  like SimpleITK or NiBabel cite\. This the dependency chain maintains a layer of indirection that can lead to memory management issues in large-scale training \cite{perezgarcia2021torchio}.

In contrast, MedImages.jl leverages Julia's type system to bind metadata directly within the \texttt{MedImage} struct. As summarized in Table \ref{tab:metadata_comparison}, by removing the runtime overhead associated with dictionary-based metadata access, this approach contributes to the improved performance observed in our benchmarks. Furthermore, the native implementation of core geometric transformations in Julia enables MedImages.jl to maintain the automatic differentiation graph across processing steps. This is particularly relevant for imaging workflows that incorporate scientific machine learning, where gradient information may be lost when operations are passed to non-differentiable compiled backends cite\. 

\begin{table}[H]
\centering
\caption{Comparative Metadata Architectures in Modern Frameworks}
\begin{tabular}{|p{2.5cm}|p{3.5cm}|p{3.5cm}|p{3.5cm}|}
\hline
\textbf{Framework} & \textbf{Metadata Storage Mechanism} & \textbf{Preservation Strategy} & \textbf{Potential Interoperability Constraints} \\
\hline
SimpleITK & Internal C++ Struct & Encapsulated in \texttt{itk::Image} & User-dependent retention (often discarded upon conversion to NumPy). \\
\hline
MONAI & MetaTensor (Dict-based) & Metadata dictionary moves with tensor & Dictionary lookup overhead; potential orientation flip during ITK interoperability. \\
\hline
TorchIO & Subject Class Attributes & Multi-modal synchronization & Potential memory buildup during extensive transformation histories. \\
\hline
MedImages.jl & Typed Julia Struct & Explicit binding within struct & N/A (Native Julia stack). \\
\hline
\end{tabular}
\label{tab:metadata_comparison}
\end{table}

\subsection*{The Primacy of Clinical and Temporal Metadata}
The lack of standardization in medical imaging preprocessing pipelines can pose a reproducibility challenge \cite{Esteban_2018}. In such settings, metadata drift, defined as the decoupling of spatial or clinical metadata from voxel data, may introduce clinically relevant errors. For example, it can lead to displacement of treatment margins in radiotherapy or compromise quantitative interpretation in PET/CT when decay-related timing information is lost \textbackslash{}cite. Quantitative nuclear medicine is particularly dependent on the integrity of clinical and temporal metadata, as standardized uptake value (SUV) calculations explicitly require accurate injection and acquisition times. In theranostic workflows, this requirement must be maintained while applying identical geometric transformations across multiple co-registered modalities, such as [$^{177}$Lu]Lu-PSMA SPECT AC/NAC, absorbed-dose maps, and CT.

MedImages.jl addresses this by design: the \texttt{BatchedMedImage} allows single-call, batched geometric transforms to be applied consistently across PET and CT modalities without manual bookkeeping. Our cross-platform validation confirms that the SUV conversion factors extracted by MedImages.jl achieve mathematical identity with the established SimpleITK benchmarks to a precision of $10^{-14}$. This numerical equivalence ensures that researchers transitioning to Julia can maintain absolute continuity with legacy clinical findings while gaining the performance required for massive biobank-scale analysis.

As demonstrated in the last experiment in these series, the batched data format shows the ability to apply consistent geometric transformations across multiple functional (SPECT, Dosemap) and anatomical (CT) data in a single operation. It ensures precise pixel-wise spatial alignment, which is essential for image augmentation and multimodal registration tasks, e.g. for maintaining consistent dosimetric outcomes.  

\subsection*{Addressing Challenge 3 (Differentiability): Navigating the Framework Ecosystem}

The Julia ecosystem is uniquely positioned for SciML because its automatic differentiation frameworks (Zygote.jl, Enzyme.jl) can differentiate through arbitrary native code, including custom spatial transformations. Nuclear medicine imaging is inherently quantitative, but its key readouts (e.g., SUV, time–activity curves, absorbed dose) are highly sensitive to geometric preprocessing such as resampling, rotation, and registration. In practice, these operations must often be integrated into AI pipelines for motion correction, longitudinal alignment, and learned preprocessing, where transform parameters are optimized jointly with network weights. However, many commonly used imaging backends execute spatial resampling outside the automatic differentiation graph, preventing end-to-end training through geometric steps \cite{Berman_2024}. To demonstrate that MedImages.jl enables fully differentiable volumetric spatial operations suitable for nuclear medicine workflows, we implement several SciML validation tests.

\paragraph{When established Python frameworks remain the pragmatic choice.}
For teams with substantial legacy Python codebases, switching frameworks carries non-trivial migration costs. MONAI provides a mature library of ready-made components---pretrained segmentation networks, standardized data-loading pipelines, and extensive augmentation transforms---that can be deployed with minimal custom code. Similarly, SimpleITK offers a comprehensive catalogue of classical filters whose correctness has been validated across decades. When the research question can be answered with an existing pipeline and the primary requirement is rapid prototyping within a specific tensor framework (like PyTorch), these tools remain highly efficient.

MedImages.jl offers a path forward---a future where imaging physics is integrated directly into the deep learning pipeline, where hardware acceleration is transparent and ubiquitous, and where the metadata that defines our physical world is preserved with the same rigor as the pixel values themselves. By bridging the divide between the physics of acquisition and the logic of inference, and by addressing the entwined challenges of \textbf{Volume} and \textbf{Speed} with a 7.2$\times$ faster turnaround than traditional caching frameworks \cite{Shi_2023}, MedImages.jl empowers researchers and clinicians to unlock the full potential of medical imaging as a precise, quantitative science. This performance is especially critical for large-scale multi-center studies and meta-analyses, where the ability to process thousands of subjects with absolute metadata integrity enables research that was previously computationally infeasible.


\subsection*{Limitations}
While MedImages.jl offers significant advantages in execution speed and differentiability, adopting the Julia ecosystem presents certain trade-offs. The most notable is the "time-to-first-plot" or compilation latency. Because Julia uses Just-In-Time compilation, the first execution of a complex spatial transformation incurs a compilation penalty, making initial startup slower than equivalent interpreted Python code. Furthermore, running full test suites or deploying models with heavy dependencies (e.g., CUDA.jl, Lux.jl, Zygote.jl) requires significant development memory, occasionally causing out-of-memory (OOM) issues on constrained hardware. Finally, while the Julia machine learning ecosystem is growing rapidly, it does not yet match the sheer volume of pre-trained models and plug-and-play architectural components available in the PyTorch/MONAI ecosystem.


\subsection{Limitations and Simplifications of UDE dosimetry calculation}
The primary objective of this experiment was to serve as a proof of concept, demonstrating the validity and feasibility of the UDE-based approach in Julia. While the UDE model definitively achieves the highest accuracy among the tested baselines, several physical abstractions were employed to simplify the initial implementation. To maintain scientific rigor and satisfy future clinical peer review, we acknowledge the following physical realities that the current framework abstracts for example :
1. \textbf{Static Anatomy}: The model assumes a rigid day-0 CT geometry, neglecting 4D respiratory motion.
2. \textbf{Isotropic PK}: Diffusion between neighboring voxels is currently ignored in the interstitial compartment.
3. \textbf{Receptor Dynamics}: Receptor capacity ($B_{MAX}$) is treated as a constant, neglecting radiation-induced receptor degradation (``stunning'').
4. \textbf{Transport Lumping}: Local $\beta^-$ and penetrating $\gamma$ transport are lumped into a single neural correction term rather than a dual-kernel separation.

Despite these necessary simplifications, the experiment successfully proves the validity of the approach: intertwining known physics with deep learning via Julia's SciML stack yields substantially superior dosimetry results compared to pure deep learning models and classical analytical convolutions.


\section{Conclusions}
The development of MedImages.jl addresses core computational and metadata challenges in modern medical imaging. By leveraging Julia's compilation architecture and type system, the library provides a unified framework where physical accuracy, computational performance, and end-to-end differentiability are achieved simultaneously. The native implementation of spatial transformations allows for transparent integration with SciML pipelines, offering a robust foundation for the next generation of physics-informed models in molecular imaging and quantitative theranostics.
In practice, we envision a complementary workflow: established MONAI/SimpleITK pipelines for routine clinical deployment, and MedImages.jl for research frontiers where performance, differentiability, or algorithmic novelty are the primary requirements.

\section*{Availability and Future Directions}
MedImages.jl is open-source and available at \url{https://github.com/JuliaHealth/MedImages.jl}. All experiments described in this paper, including performance benchmarks and differentiability validations, are reproducible via the scripts provided in the repository's \texttt{experiment} directory.

Future work will focus on:
\begin{enumerate}
    \item Integrating with advanced reconstruction frameworks (KomaMRI.jl \cite{CastilloPassiCoronado2023}, GIRFReco.jl \cite{JaffrayWu2024}) for end-to-end physics-informed learning.
    \item Implementing differentiable registration algorithms leveraging Zygote.jl.
    \item Expanding support for standardized formats like DICOM-SEG via Highdicom \cite{BridgeGorman2022} integration.
    \item Developing architectural foundations for clinical-grade AI agents (e.g., VoxelPrompt) capable of multi-modal "Medical Superintelligence" benchmarks where absolute metadata integrity is a prerequisite for safe automated diagnostics.
\end{enumerate}

\section*{Declarations}
\textbf{Ethics approval and consent to participate:} This research relies solely on retrospective analysis of fully anonymized, publicly available datasets. No new human subjects data was collected. Therefore, formal ethics committee approval and informed consent were not required. \\
\textbf{Declaration of competing interests:} The authors declare that they have no known competing financial interests or personal relationships that could have appeared to influence the work reported in this paper. \\
\textbf{Funding sources:} This research did not receive any specific grant from funding agencies in the public, commercial, or not-for-profit sectors. \\
\textbf{Declaration of generative AI and AI-assisted technologies:} During the preparation of this work, the author(s) used generative AI and AI-assisted technologies for improving style and assisting in code development. After using these tools, the authors reviewed and edited the content as needed and take full responsibility for the content of the published article.





\section{Appendix}

Architectural Foundations: The Coordinate System and Metadata
To accurately integrate functional data with anatomical data across different spatial resolutions, software must rigorously maintain the mapping between discrete voxel indices and physical patient space. This mapping from voxel coordinates $\mathbf{v} = [i, j, k, 1]^\top$ to continuous world coordinates $\mathbf{p} = [x, y, z, 1]^\top$ relies on a $4 \times 4$ homogeneous affine transformation matrix $M$:
\begin{equation}
\mathbf{p} = M \mathbf{v}, \quad M = 
\begin{bmatrix} 
R S & T \\
\mathbf{0} & 1
\end{bmatrix}
\end{equation}
where $R$ is the $3 \times 3$ rotation matrix, $S$ is a diagonal scaling matrix (voxel spacing), and $T$ is the translation vector (origin). MedImages.jl ensures that any spatial operation $\mathcal{T}$ applied to the voxel grid is mathematically coupled with an update $M_{new} = M_{old} \times M_{\mathcal{T}}$, preventing spatial desynchronization.

Furthermore, quantitative nuclear medicine relies on precise temporal and clinical metadata. The Standardized Uptake Value (SUV), a critical metric for PET/CT, is calculated as:
\begin{equation}
SUV = \frac{C(t)}{ID / BW}
\end{equation}
where $C(t)$ is the decay-corrected radioactive concentration, $ID$ is the injected dose, and $BW$ is patient body weight. Discarding parameters such as \textit{Acquisition Time} or \textit{Patient Weight} compromises the quantitative utility of the image. MedImages.jl intrinsically protects these fields within its data structures.

\subsection{Workflow}


\paragraph{When MedImages.jl offers a clear advantage for SciML.}
The Julia-based approach becomes compelling when research transcends standard convolutional network training and enters the realm of physics-informed models or custom numerical workflows \cite{RoeschGreener2023}:
\begin{itemize}
    \item \textbf{Unifying the Fractured SciML Landscape.} In Python, moving a differentiable model from PyTorch to a JAX-based differential equation solver often requires significant code rewriting or bridging overhead. Julia's multiple dispatch solves the expression problem, meaning essentially all packages are interoperable natively \cite{PalBhattacharya2024, EschleGl2023}. MedImages.jl serves as the critical entry point to make this uniquely unified SciML ecosystem accessible for medical imaging.
    \item \textbf{Novel algorithm development.} When researchers need to implement a new spatial transformation, loss function, or physics-informed model from scratch, Julia's single-language stack eliminates the need to write C++/CUDA extensions. The algorithm can be prototyped, debugged, and GPU-accelerated within the same language, with full automatic differentiation support.
    \item \textbf{Cross-domain scientific computing.} Julia's ecosystem uniquely bridges medical imaging with adjacent scientific domains---DifferentialEquations.jl for pharmacokinetic modelling, KomaMRI.jl \cite{CastilloPassiCoronado2023} for MRI pulse-sequence simulation, Flux.jl/Lux.jl \cite{pal2023lux} for deep learning---all within a single runtime. This composability enables workflows that would otherwise require fragile inter-process communication between Python, MATLAB, and C++ components.
    \item \textbf{Differentiable imaging physics.} Tasks such as learned registration, differentiable augmentation, or physics-informed neural networks require gradients to flow through spatial transformations. MedImages.jl's native Julia kernels integrate directly with Zygote.jl and Enzyme.jl, whereas MONAI and TorchIO must fall back on non-differentiable C++ backends for many geometric operations.
\end{itemize}

We do not advocate replacing established tools wholesale. Rather, we recommend choosing the framework that best matches the task at hand, especially when differentiability and complex scientific computing are core requirements.

\subsubsection*{Code Example}
The following snippet demonstrates loading a multimodal pair, rotating them simultaneously using the batched structure, and computing the SUV:

\begin{lstlisting}[basicstyle=\ttfamily\footnotesize]
using MedImages

# Load PET and CT (DICOM metadata is preserved)
pet = load_image("patient_1/pet.dcm")
ct = load_image("patient_1/ct.nii")

# Create a batched tensor to process simultaneously
batch = create_batched_medimage([pet, ct])

# Rotate both volumes by 45 degrees around Z-axis (GPU supported)
rotated_batch = rotate_mi(batch, 45.0)

# Extract SUV statistics for the PET scan using an ROI mask
suv_stats = calculate_suv_statistics(rotated_batch[1], roi_mask)
println("Mean SUV: ", suv_stats.mean_suv)
\end{lstlisting}

\begin{table}[H]
\centering
\caption{Detailed information  - Experimental Series 2: the batched multi-modal theranostic processing}
\label{tab:method_info_batched_Lu177}
\begin{tabular}{|p{0.68\textwidth}|p{0.24\textwidth}|}
\hline
\textbf{Methodology}& \textbf{Study-specific information} \\
\hline
Number of patients / datasets used (single case vs.\ multiple cases) & \\
\hline
Source and acquisition details of the imaging data (scanner type, reconstruction parameters) & \\
\hline
Exact voxel dimensions and spatial resolution of each modality before transformation & \\
\hline
Whether modalities were pre-registered or inherently aligned before batching & \\
\hline
Definition of the rotation (angle, axis, interpolation method) & \\
\hline
Interpolation schemes used per modality (e.g., linear, nearest-neighbor) & \\
\hline
Handling of differing voxel grids during batched transformation (internal resampling strategy) & \\
\hline
Quantitative evaluation metrics, if any, beyond visual inspection & \\
\hline
Software and hardware environment (CPU/GPU, Julia version) & \\
\hline
Reproducibility details (random seeds, code availability, parameter settings) & \\
\hline
\end{tabular}
\end{table}
```
\end{document}

